% This Source Code Form is subject to the terms of the Mozilla Public
% License, v. 2.0. If a copy of the MPL was not distributed with this
% file, You can obtain one at https://mozilla.org/MPL/2.0/.
\documentclass[11pt,a4paper]{article}

\usepackage[margin=2.2cm]{geometry}
\usepackage{amsmath,amssymb,amsthm}
\usepackage{booktabs}
\usepackage{array}
\usepackage{enumitem}
\usepackage{needspace}
\usepackage{microtype}
\usepackage[hidelinks]{hyperref}

% The same source gives the same PDF: no dates, no random identifier and no
% description of the machine that built it. `just pdfs` passes the SHA-256
% of this file as \sourcehash, and the PDF records it, so that
% `just pdfs-check` can tell whether the committed PDF was built from the
% current source.
\ifdefined\pdfinfoomitdate
  \pdfinfoomitdate=1
  \pdftrailerid{}
  \pdfsuppressptexinfo=-1
\fi
\providecommand{\sourcehash}{unknown}
\hypersetup{
  pdftitle={Optimal playback with exchanges between the sender and each receiver},
  pdfsubject={Built from docs/star-optimum.tex with SHA-256 \sourcehash}
}

\setlist{itemsep=2pt,topsep=4pt,parsep=0pt}
\emergencystretch=3em
\clubpenalty=10000
\widowpenalty=10000

\newcommand{\lean}[1]{\texttt{#1}}
\newcommand{\R}{\mathbb{R}}
\newcommand{\E}{\mathbb{E}}

\theoremstyle{definition}
\newtheorem{definition}{Definition}[section]
\newtheorem{theorem}[definition]{Theorem}
\newtheorem{remark}[definition]{Remark}

\begin{document}

\begin{center}
{\LARGE\bfseries Optimal playback with exchanges between\\[2pt]
the sender and each receiver\par}
\end{center}
\medskip

This document presents a solution to the problem of
\path{docs/specification.pdf} in which the hosts form a \emph{star}: every
message goes between the sender and one \emph{receiving host}, a host that
drives a speaker and is not the sender's. It proves that the solution's way
of playing is the best possible for the quadratic objectives, among all ways
that use only what a receiving host has seen, and compares it with the
baseline of \path{docs/fixed-playout-delay.pdf}. The proofs are in
\path{lean/Unisono/Solutions/Star/}, the model of the world is the
second-order model of the specification
(\path{lean/Unisono/Spec/Models/SecondOrder.lean}), and the simulation is in
\path{crates/unisono-sim/src/star.rs}. The numbers below come from
\texttt{cargo run --release -p unisono-sim --example star}.

The solution has three parts:
\begin{itemize}
\item a \emph{front end} that keeps, of each group of consecutive clock
  exchanges, the one with the shortest round trip;
\item an \emph{estimator} that computes the best estimate of the offset and
  the relative frequency of the sender's clock and the host's clock from the
  front end's measurements: a Kalman filter, which in the stationary
  Gaussian case is the Wiener filter;
\item the \emph{law} of \path{lean/Unisono/Model/Regulation.lean}, which
  steers the playback position towards the estimate with the least squared
  distance and the least squared change of pace.
\end{itemize}
The main theorem says that the law applied to the best estimate, the
conditional mean, is optimal among all causal ways of playing, for any
probability distribution of the noise. With several receivers, a
broadcast from the sender removes the sender's clock from the
synchronization error. The latency level comes from the distribution of the
holds.

\section{The problem of one receiver}

A receiver reads the stream once in each cycle of its device, the callback
in which the device asks for the next block. Let:
\begin{itemize}
\item $S_k$ be the position, in frames, that is truly due at the start of
  cycle $k$: the position whose sound was captured the latency $L$ before
  the cycle's first output sample leaves;
\item $p_k$ be the receiver's position at the start of cycle $k$, and $n_k$
  the number of frames that cycle $k$ reads, so that $p_{k+1} = p_k + n_k$;
\item $w_k = n_k - n_{k-1}$ be the change of the pace at cycle $k$;
\item $\mathcal F_k$ be the $\sigma$-algebra of what the receiver has seen
  before cycle $k$.
\end{itemize}
A way of playing is \emph{causal} when $p_0$, $n_{-1}$ and every $n_k$ are
$\mathcal F_k$-measurable with finite mean squares. For a gain
$0 < g < 1$ and the weight
\[
  c_g = \frac{g^4}{(2 - g)^2 (1 - g)},
\]
which takes every positive value once as $g$ ranges over $(0, 1)$, cycle $k$
costs
\[
  c_g\,(S_k - p_k)^2 + w_k^2 .
\]
The first term is the squared latency error at the start of the cycle,
counted in frames: the error in seconds times the nominal rate. The second
term is the squared change of the playback speed between two cycles, times
the square of the cycle's length in frames. The specification's objectives
count every output sample; this document counts the first output sample of
each cycle.

\subsection{The orthogonality principle}

In a real inner product space, let $K$ be a linear subspace, $b$ a vector
and $x^\circ \in K$ a vector with $b - x^\circ$ orthogonal to $K$. Then
every $x \in K$ satisfies
$\|b - x\|^2 = \|b - x^\circ\|^2 + \|x - x^\circ\|^2$. The vector
$x^\circ$ is the orthogonal projection of $b$ on $K$, and no element of
$K$ is closer to $b$. For random variables with finite mean squares, the
mean of the product is an inner product, and the $\mathcal F_k$-measurable
random variables form a closed subspace. The conditional expectation
$\E[S \mid \mathcal F_k]$ is the orthogonal projection on that subspace,
which Mathlib provides.

The principle is used twice. First, the receiver's position $p_k$ is
$\mathcal F_k$-measurable, so the error of the best estimate
$\hat S_k = \E[S_k \mid \mathcal F_k]$ is orthogonal to the distance
$\hat S_k - p_k$, and
\[
  \E(S_k - p_k)^2 = \E(\hat S_k - p_k)^2 + \E(S_k - \hat S_k)^2
\]
(\lean{mean\_square\_split}). Second, the revisions of the best estimates
between two cycles are orthogonal to everything known before the first of
them, and against such revisions the law is optimal
(\lean{mean\_cost\_sum}).

\subsection{The best estimate and the law are optimal together}

Let $\hat a_k = \E[S_{k+1} - S_k \mid \mathcal F_k]$ be the expected
advance over cycle $k$. The \emph{state} before cycle $k$ consists of the
distance $E_k = \hat S_k - p_k$ of the position behind the best estimate,
and the number $v_k = n_{k-1} - \hat a_k$ of frames by which the last cycle
read more than the expected advance. The \emph{revision} between cycle $k$
and cycle $k + 1$ is
\[
  \xi_k = \bigl(\E[S_{k+1} \mid \mathcal F_{k+1}] - \E[S_{k+1} \mid
  \mathcal F_k],\ \hat a_k - \hat a_{k+1}\bigr).
\]
The law with gain $g$ changes the pace by
\[
  \ell_g(E, v) = \frac{g^2}{2 - g}\,E - g\,v ,
\]
and the cost of all cycles from a state under the law, without revisions,
is the quadratic form
\[
  V_g(E, v) = \frac{g^3}{(2 - g)(1 - g)}\,E^2 - \frac{2 g^2}{2 - g}\,E v
    + g\,v^2
\]
(\lean{value}).

\begin{theorem}[The cost splits]\label{thm:split}
Suppose the pace of the due position is expected to stay:
$\E[S_{k+2} - S_{k+1} \mid \mathcal F_k] = \E[S_{k+1} - S_k \mid
\mathcal F_k]$ for every $k$, and every $S_k$ has a finite mean square. Then
for every causal way of playing and every number $n$ of cycles,
\begin{multline*}
  \sum_{k < n} \E\bigl[c_g (S_k - p_k)^2 + w_k^2\bigr] + \E V_g(E_n, v_n)
  = c_g \sum_{k < n} \E(S_k - \hat S_k)^2 + \E V_g(E_0, v_0) \\
  + \sum_{k < n} \E V_g(\xi_k)
  + \frac{1}{1 - g} \sum_{k < n} \E\bigl(w_k - \ell_g(E_k, v_k)\bigr)^2
\end{multline*}
(\lean{mean\_cost\_split}).
\end{theorem}

No way of playing changes the first three terms on the right. The last term
is zero exactly when every change of the pace is the law's change applied to
the state. Playing that way is causal (\lean{lawRead}, \lean{known\_lawRead})
and therefore optimal: for every causal way of playing from the same start,
its cost plus the value of the final state is at least the optimal one
(\lean{mean\_cost\_law\_le}). The theorem asks nothing of the distribution
of the noise: the revisions need not be independent, of one law, or
Gaussian. The only condition is that the pace is expected to stay, which a
due position driven by clocks with white and random-walk frequency noise
satisfies. Flicker frequency noise does not satisfy it, because its
frequency is not a martingale, and is outside the theorem.

Together with the law's long-run results (\lean{mean\_cost\_rate},
\lean{mean\_value\_law\_le}), the theorem gives the long-run statement: the
mean cost per cycle of any causal way of playing is at least $c_g$ times the
mean squared error of the best estimate plus the mean value of a revision,
and the law attains this bound, because its states keep a bounded mean
value. This combination is not stated as one theorem in Lean.

The theorem holds for every filtration. It therefore applies to whatever the
receiver chooses to observe, including the output of the front end. Which
estimator gives the conditional mean depends on the model. When the front
end's errors and the clocks' noise are Gaussian, the conditional mean is the
Kalman filter, and the theorem makes the Kalman filter with the law optimal
among all causal ways of playing. For other distributions, the Kalman filter
is the best linear estimator; the corresponding statement for linear ways of
playing follows from the same argument with linear projections, but is not
formalized.

\subsection{The frequency domain}

The same solution comes out of the Wiener--Hopf method. In the stationary
Gaussian case with a measurement at every cycle, the optimal map from the
measurements to the positions is linear and time-invariant, and the
Wiener--Hopf method finds it by factoring two spectra.

\paragraph{The weight.} The change of the pace is the second difference of
the position. On the unit circle $z = e^{i\omega}$, the cost's weight is
therefore $c_g + |1 - z|^4$. It factors as
\[
  c_g + \bigl(2 - z - z^{-1}\bigr)^2 = \frac{P_g(z)\,P_g(z^{-1})}{1 - g},
  \qquad P_g(z) = z^2 - \frac{4(1 - g)}{2 - g}\,z + (1 - g)
\]
(\lean{weight\_factor}), where $2 - z - z^{-1} = |1 - z|^2$ on the circle.
Under the law, two cycles without revisions satisfy
$x_{k+2} = \frac{4(1 - g)}{2 - g}\,x_{k+1} - (1 - g)\,x_k$
(\lean{step\_step}): the law's closed loop has the characteristic
polynomial $P_g$. Both roots of $P_g$ lie inside the unit circle, with
modulus $\sqrt{1 - g}$. The law is thus the causal factor of the weight,
as the Wiener--Hopf solution prescribes.

\paragraph{The papers' law.} The earlier papers derived, in continuous time,
the law that minimizes $q e^2 + r \dot u^2$ for a latency error $e$ driven by
white noise, where $u$ is the warp: the difference between the playback
speed and the nominal speed. The weight $q + r\omega^4$ factors as
$|W(i\omega)|^2$ with $W(s) = c_2 s^2 + c_1 s + c_0$, $c_0^2 = q$,
$c_2^2 = r$ and $c_1^2 = 2 c_0 c_2$ (\lean{butterworth}): the
second-order Butterworth polynomial with natural frequency
$(q/r)^{1/4}$ and damping $1/\sqrt 2$, which is the papers' closed loop.
The law of this document is its exact discrete-time counterpart: when the
measurement is exact and the due position is a random walk, the estimator
is the measurement itself, and the optimal controller is the causal factor
$P_g$ of the discrete weight.

\paragraph{The estimate.} Let the due position follow a random walk whose
steps have variance $q$, measured at every cycle with white noise of
variance $r$. Let $a$ be the root in $(0, 1)$ of
$a^2 - (2 + q/r)\,a + 1 = 0$ (\lean{filterRoot}). The spectrum of the
measurements factors as
\[
  \frac{q + r\,(2 - z - z^{-1})}{(1 - z)(1 - z^{-1})}
  = \frac{r}{a}\,\frac{(1 - a z)(1 - a z^{-1})}{(1 - z)(1 - z^{-1})}
\]
(\lean{estimate\_factor}), and the Wiener filter is
$\hat S = \frac{1 - a}{1 - a z^{-1}}\,y$ for measurements $y$. In the time
domain this filter is the steady state of the Kalman filter, whose gain is
$1 - a$ (\lean{riccati\_fixed}). Its estimate errs with variance
$r(1 - a)$ (\lean{posterior\_variance}), and its revisions have variance
$q$ (\lean{revision\_variance}). The optimal mean cost per cycle is
therefore
\[
  c_g\, r\,(1 - a) + \frac{g^3}{(2 - g)(1 - g)}\,q ,
\]
the first term from the estimate's error and the second from the value of a
revision, which moves only the position. The simulation's Kalman filter
reaches the variance $r(1 - a)$ to nine digits
(\texttt{the\_filter\_reaches\_the\_closed\_form}).

The clocks of the second-order model add random-walk frequency noise, which
gives the estimator a second state, the relative frequency. Its steady
state solves a $2 \times 2$ Riccati equation, which the simulation iterates.
Flicker frequency noise has no finite-dimensional state; its spectrum
$h_{-1}/f$ is not rational, and its spectral factor is computed from the
logarithm of the spectrum by Fourier transforms. Neither case is formalized
here.

\subsection{The front end}

Network delays are one-sided: a message never arrives before the link's
smallest delay, and holds only make messages later. A linear filter of $n$
samples errs by about $\sigma/\sqrt n$, where $\sigma$ is the spread of the
delays. The shortest of $n$ delays errs by about $1/(n f_0)$, where $f_0$
is the density of the delays just above the smallest. Of each group of
consecutive exchanges, the front end therefore keeps the one with the
shortest round trip, and the estimator works on these measurements. A hold
that does not cover a whole group does not reach the measurement. What
remains is the asymmetry of the two directions' smallest delays, which no
exchange with the sender can observe.

\section{Several receivers}

Receiving host $j$ must follow $\theta_j = c + i_j$, the difference between
the sender's clock and its own. The part $c$ is common to all receiving
hosts: the sender's clock, and the delays on the sender's side. The part
$i_j$ belongs to host $j$: its clock and its link. Host $j$ measures
$\theta_j + m_j$, where $m_j$ is the error of its measurements. A linear way
of playing applies a linear map $F_j$ to the measurements, such as a causal
filter acting on sequences of random variables, and makes the error
$e_j = F_j(\theta_j + m_j) - \theta_j$. Synchronization depends only on the
differences $e_j - e_{j'}$.

\begin{theorem}[What the receivers share cancels]\label{thm:cancel}
\leavevmode
\begin{itemize}
\item If all receivers apply the same map $F$, then
  $e_j - e_{j'} = F(i_j + m_j - i_{j'} - m_{j'}) - (i_j - i_{j'})$, and
  the synchronization cost does not depend on $c$
  (\lean{error\_sub\_error\_same}, \lean{syncCost\_same\_filter}).
\item If two receivers apply different maps, the common part enters their
  difference as $(F_j - F_{j'})\,c$ (\lean{error\_sub\_error}).
\item Suppose the sender broadcasts an estimate $c + \eta$ of the common
  part. Each receiver applies a map $G$, the same for all, to the broadcast,
  and its own map $F_j$ to its measurement minus the broadcast. Then
  $e_j - e_{j'} = F_j(i_j + m_j - \eta) - F_{j'}(i_{j'} + m_{j'} - \eta)
  - (i_j - i_{j'})$, whatever the maps $F_j$ are, and the synchronization
  cost does not depend on $c$ (\lean{broadcastError\_sub},
  \lean{syncCost\_broadcast}).
\end{itemize}
\end{theorem}

The identities hold for every realization, so they hold for the mean
squares as well. The first explains why identical receivers do not need the
broadcast. The second explains why a mix of wired and wireless receivers
does: the optimal filter of a receiver with noisy measurements is narrower,
so it follows the sender's clock with a different lag. The broadcast's own
error $\eta$ still enters through the receivers' maps. The sender averages
the receivers' measurements with weights inversely proportional to their
variances, so $\eta$ is mostly the error of the best-measured receivers.

\paragraph{Numbers.} The simulation of the second-order model in
Table~\ref{tab:broadcast} has three wired receivers, whose measurements err
by $3\,\mu$s, and two wireless receivers, whose measurements err by
$150\,\mu$s, one measurement every two seconds. All receivers' clocks are
quartz: white frequency noise of $2 \times 10^{-18}$ per hertz, an Allan
deviation of $10^{-9}$ at one second, and a random walk of the frequency
that adds $10^{-18}$ to its variance per second. The sender's clock is
either such a quartz, or one whose frequency wanders a hundred times faster,
as an oscillator that heats and cools does. Without the broadcast, each
receiver runs its own Kalman filter. With it, every receiver runs the same
filter on the broadcast and its own filter on the rest.

\begin{table}[ht]
\centering
\begin{tabular}{lrrrrrr}
\toprule
& \multicolumn{2}{c}{all receivers} & \multicolumn{2}{c}{wired pairs}
  & \multicolumn{2}{c}{wired and wireless} \\
sender's clock & own & broadcast & own & broadcast & own & broadcast \\
\midrule
quartz & 7.05 & 6.66 & 0.89 & 0.87 & 12.70 & 11.89 \\
wandering & 11.52 & 6.66 & 1.33 & 0.87 & 21.03 & 11.89 \\
\bottomrule
\end{tabular}
\caption{Root mean square spread of the errors, in microseconds: the
standard deviation over all receivers, and the difference within the pairs
of each kind.}\label{tab:broadcast}
\end{table}

With a quartz sender the broadcast gains little: the receivers' lags differ
by less than their own noise. With a wandering sender, the receivers without
the broadcast drift apart, while with the broadcast the spreads are the same
as with a quartz sender. The sender's clock no longer enters them, as
Theorem~\ref{thm:cancel} says.

\paragraph{What a star cannot do.} A change in the asymmetry of a link's two
directions shifts the offset estimate by half the change, and no exchange
with the sender reveals it. On a wireless link, a change of the link's rate
moves the smallest delay of each direction separately. Exchanges between the
receivers, or the arrival times of one multicast message at several
receivers, would expose such changes. They are left to a solution that
lets receivers communicate with each other.

\section{The latency level from the holds}

A hold is an episode in which a link keeps the messages sent on it and
releases them when it ends. Suppose the latency level leaves a margin of
$m$ seconds above what the other delays use. Then a hold of length $h$
makes the receiver miss the output samples of $(h - m)^+$ seconds, where
$x^+$ is the larger of $x$ and zero. When holds start at the rate $\lambda$
per second with lengths distributed as $\mu$, the mean time missed per
second is $\lambda \int (h - m)^+ \,d\mu(h)$, by Campbell's formula for
stationary marked point processes, as long as holds do not overlap and
playback resumes on time after each of them. With a price $a$ per second of
latency and a price $b$ per second of missed playback, the margin $m$ costs
\[
  a\,m + b \int (h - m)^+\, d\nu(h), \qquad \nu = \lambda\,\mu .
\]
When several receivers play at one common level, their measures add, each
shifted by what the other delays use at that receiver.

\begin{theorem}[The best margin is a quantile]\label{thm:margin}
Let $\nu$ be a finite measure with a finite mean, and let
$b\,\nu\{h > m\} \le a \le b\,\nu\{h \ge m\}$. Then every margin $x$ costs
at least as much as $m$ (\lean{margin\_best}).
\end{theorem}

The best margin is the one beyond which holds are rarer than $a/b$ per
second; only the tail of their distribution matters. For holds of $30$~ms
on average, exponentially distributed, every $20$ seconds,
$\lambda e^{-m/\bar h} = a/b$ gives $m = \bar h \ln(\lambda b / a)$, where
$\bar h$ is the mean length. Table~\ref{tab:holds} shows the margins for
three ratios of the prices, the fraction of time that Campbell's formula
predicts to be missed, $\lambda \bar h e^{-m/\bar h}$, and the fraction of
packets that the simulated link makes late by more than the margin, out of
$2 \times 10^7$ packets.

\begin{table}[ht]
\centering
\begin{tabular}{rrrr}
\toprule
$a/b$ & margin (ms) & Campbell & simulated \\
\midrule
$10^{-2}$ & 48.3 & $3.00 \times 10^{-4}$ & $2.88 \times 10^{-4}$ \\
$10^{-3}$ & 117.4 & $3.00 \times 10^{-5}$ & $3.02 \times 10^{-5}$ \\
$10^{-4}$ & 186.4 & $3.00 \times 10^{-6}$ & $2.80 \times 10^{-6}$ \\
\bottomrule
\end{tabular}
\caption{Margins from the quantile rule for holds of $30$~ms on average every
$20$~s, and the fraction of time missed.}\label{tab:holds}
\end{table}

The baseline needs a latency above the longest delay, holds included, to
guarantee that it misses nothing. The quantile rule instead prices each
millisecond of margin against the playback it saves.

\section{Against the baseline}

Table~\ref{tab:baseline} compares the baseline with the solution of this
document in the simulation of \path{crates/unisono-sim/src/star.rs}. Each
run lasts ten minutes, with cycles of $10$~ms, and the first minute is not
counted. Every receiver sends a clock query every $0.25$~s, so both ways of
playing use the network equally. Excess delays are exponentially
distributed. The wired links add $5\,\mu$s on average to nine messages in
ten and $200\,\mu$s on average to the tenth. The wireless links add
$0.3$~ms on average to half of the messages and $3$~ms on average to the
other half, and hold everything for $30$~ms on average every $20$ seconds.
The clocks are quartz,
with relative frequencies of up to $10$ parts per million. The baseline
steers towards the estimate of the latest exchange with a slew of
$10^{-3}$ seconds per second. The solution keeps the exchange with the
shortest round trip out of every eight, filters with the variance of its
link's measurements, and uses the law with gain $g = 0.05$, whose closed
loop settles in about $0.4$~s.

\begin{table}[ht]
\centering
\small
\setlength{\tabcolsep}{3.5pt}
\begin{tabular}{lrrrrrrrr}
\toprule
& \multicolumn{2}{c}{spread ($\mu$s)} & \multicolumn{2}{c}{error ($\mu$s)}
  & \multicolumn{2}{c}{speed change (ppm)}
  & \multicolumn{2}{c}{largest change (ppm)} \\
receivers & base & solution & base & solution & base & solution & base
  & solution \\
\midrule
two wired & 31.7 & 0.24 & 45.2 & 0.42 & 186.8 & 0.002 & 2000 & 0.04 \\
two wireless & 226.7 & 21.1 & 321.9 & 24.2 & 292.0 & 0.20 & 2000 & 6.5 \\
three wired, two wireless & 186.1 & 14.6 & 210.6 & 17.3 & 234.2 & 0.11
  & 2000 & 5.1 \\
\bottomrule
\end{tabular}
\caption{The baseline against the solution. The spread is the root mean
square of the standard deviation of the receivers' errors; the error is the
root mean square of each receiver's error against the sender's clock; the
speed change is the change of the playback speed from one cycle to the
next, in parts per million.}\label{tab:baseline}
\end{table}

On the wired links the solution synchronizes a hundred times more closely.
It changes the speed by at most a few hundredths of a part per million, where
the baseline reverses its correction at the full slew. On the wireless links the
solution's error is set by the asymmetric medium access, which the shortest
round trip of eight exchanges does not remove entirely. It is still ten
times smaller than the baseline's. The comparison simulates the clock
exchanges, not the stream: the devices run exactly on their hosts' clocks,
and the stream's late packets are simulated separately, in the previous
section.

\section{What is proved}

\begin{itemize}
\item \path{lean/Unisono/Spec/Models/SecondOrder.lean}: the second-order
  model and its non-empty example (\lean{SecondOrderModel.example\_holds} in
  \path{lean/Unisono/Spec/Checks.lean}).
\item \path{lean/Unisono/Solutions/Star/Separation.lean}: the cost splits
  for every causal way of playing (\lean{mean\_cost\_split}), and the best
  estimate with the law is optimal (\lean{mean\_cost\_law},
  \lean{mean\_cost\_law\_le}).
\item \path{lean/Unisono/Solutions/Star/Spectral.lean}: the factorization
  of the weight and the law's characteristic polynomial, the Butterworth
  factor of the papers' weight, and the steady state of the scalar Kalman
  filter.
\item \path{lean/Unisono/Solutions/Star/Receivers.lean}: the cancellation
  identities for several receivers, with and without the broadcast.
\item \path{lean/Unisono/Solutions/Star/Level.lean}: the quantile rule for
  the margin.
\end{itemize}
Every theorem depends only on the axioms \lean{propext},
\lean{Classical.choice} and \lean{Quot.sound}
(\path{lean/Unisono/Audit.lean}).

The theorems are about the problem of one receiver cycle by cycle, and
about linear maps for several receivers. The solution is not yet defined as
a solution of the specification's run, as the baseline is, so no theorem
connects it end to end with the specification's objectives. The Kalman
filter's optimality over all causal ways of playing requires Gaussian
noise, Campbell's formula is used rather than proved, and flicker frequency
noise is outside the theorem.

\end{document}
